% ---------------------------------------------------------------
% preamble.tex  (RICOSTRUITO: il file originale non era stato caricato)
% Se hai già un tuo preamble.tex, usa quello: contiene solo gli
% ambienti e le macro necessari ai file .tex (vedi sotto).
% ---------------------------------------------------------------
\usepackage[utf8]{inputenc}
\usepackage[T1]{fontenc}
\usepackage[a4paper,margin=2.5cm]{geometry}
\usepackage{amsmath,amssymb,amsthm,mathtools,mathrsfs,yhmath}
\usepackage{centernot}
\usepackage{enumitem}
\usepackage{float}
\usepackage[dvipsnames]{xcolor}
\usepackage{tikz}
\usetikzlibrary{arrows.meta,calc,decorations.markings,positioning,shapes.geometric}
\usepackage{tikz-cd}
\usepackage{caption,subcaption}
\usepackage{hyperref}

\newcommand{\Func}[5]{\[ #1\colon #2 \longrightarrow #3,\qquad #4 \longmapsto #5 \]}
\providecommand{\widearc}[1]{\wideparen{#1}}


\theoremstyle{plain}
\newtheorem{teo}{Teorema}[section]
\newtheorem{prop}[teo]{Proposizione}
\newtheorem{lema}[teo]{Lemma}
\newtheorem{coro}[teo]{Corollario}
\theoremstyle{definition}
\newtheorem{definicion}[teo]{Definizione}
\newtheorem{ejemplo}[teo]{Esempio}
\newtheorem*{ejercicio*}{Esercizio}
\theoremstyle{remark}
\newtheorem{observacion}[teo]{Osservazione}
\newtheorem{notacion}[teo]{Nota}
\newtheorem*{intuizione}{Intuizione}
